\documentclass[10pt,a4paper]{article}
\usepackage[T1]{fontenc}
\usepackage[utf8]{inputenc}
\usepackage{lmodern}
\usepackage{amsmath,amssymb}
\usepackage[margin=21mm]{geometry}
\usepackage{microtype}
\usepackage[hidelinks]{hyperref}
\usepackage{enumitem}
\setlist[itemize]{leftmargin=*,nosep,topsep=3pt}
\setlength{\parindent}{0pt}
\setlength{\parskip}{5pt}
\setlength{\emergencystretch}{3em}
\newcommand{\question}[1]{\par\medskip{\normalsize\bfseries #1}\par\smallskip}
\pagestyle{plain}

\begin{document}
{\Large\bfseries Cortical phase-amplitude coupling is key to the occurrence and treatment of freezing of gait\par}

\textbf{Paper:} Yin et al. (2022).\quad Human walking study.\par
\href{https://doi.org/10.1093/brain/awac121}{DOI: \nolinkurl{10.1093/brain/awac121}}\par
{\footnotesize Page references below use PDF page numbers, including front matter.}

\question{To which general question does this paper relate?}
Measurement, freezing of gait, and the choice of a clinical control objective. Tests whether a cortical coupling measure tracks freezing and changes with STN stimulation.

\question{How does it answer the general questions?}
\textbf{Measured and manipulated.} Motor-cortex ECoG and gait were recorded during walking tasks in 16 patients with Parkinson's disease. Beta--gamma phase-amplitude coupling (PAC) was higher in trials with freezing and related to freezing severity. STN-DBS reduced PAC and alleviated freezing (PDF pp.~1--10).

\textbf{Usable insight.} Freezing still improved with DBS when comparing trials with similar PAC. Therefore, PAC reduction alone does not explain the benefit. A proposed ``bandwidth'' account combines baseline synchrony, task demands, and processing capacity; it is a conceptual explanation rather than an identified dynamical model (PDF pp.~9--13).

\textbf{Model relevance.} The supported behavioural objective is reducing freezing while preserving walking. PAC is a candidate feedback feature, requiring a nonlinear signal-processing block beyond the notebook's linear measurement equation. The paper supplies no delayed rate-model parameterization or implemented PAC-based adaptive controller.

\question{What papers does it cite that could provide more information?}
\begin{itemize}
\item Pozzi et al. (2019), \emph{Freezing of gait in Parkinson's disease reflects a sudden derangement of locomotor network dynamics}: network context (ref.~13, PDF p.~14).
\item Cole et al. (2017), \emph{Nonsinusoidal beta oscillations reflect cortical pathophysiology in Parkinson's disease}: waveform contributions to PAC (ref.~41, PDF p.~14).
\end{itemize}

\end{document}
